\section{REFERENCIAL TEÓRICO}

\subsection{Criptografia de mensagens}

Há muitos anos, a informação é uma ferramenta crucial pra se obter
vantagens sobre seus concorrentes. Criptografá-la para impedir acesso
por terceiros, se tornou necessário para a confiabilidade da
transmissão das informações e da privacidade dos indivíduos. Assim como
\citeonline[p.~1]{carvalho2021} dizem que “Proteger mensagens e
informações sigilosas é uma prática que remota aos primórdios da
humanidade, principalmente em contexto de guerra e diplomacia”.

Ao tratarmos de comunicação por meios digitais, transportando-a por
intermédio de redes públicas de computadores, como é o caso da
internet, torna-se imprescindível garantir a segurança dessa
comunicação, de modo a garantir os fundamentos da segurança da
informação, sendo eles, disponibilidade, integridade, controle de
acesso, autenticidade, não-repudiação e privacidade, como aborda
\citeonline{oliveira2012}. Mostra-se necessária a implementação de
técnicas computacionais que visam atender os requisitos para a
proteção da comunicação. Nessa conjuntura, dois tipos de criptografia
se mostram bastante promissores: a simétrica e a assimétrica.

\subsubsection{Criptografia simétrica}

No método de chave simétrica, tanto a criptografia dos dados quanto a
descriptografia são feitas com base em uma única chave chamada de
chave privada, também podendo ser chamada de chave secreta. Um canal
seguro é necessário para compartilhar essa chave entre o remetente e
o destinatário, de modo a garantir que nenhuma pessoa não autorizada
possa obter a chave. A figura abaixo ilustra o processo de
criptografia simétrica.

\subsubsection{Criptografia assimétrica}

Já na criptografia assimétrica duas chaves são necessárias: uma é a
chave privada que deve ser mantida em segredo e a outra é a chave
pública. A criptografia é realizada usando a chave pública, já a
chave secreta é usada para descriptografar o conteúdo cifrado. Ambas
as chaves são matematicamente relacionadas entre si. É importante
ressaltar ainda, que não é possível descriptografar o texto com a
chave pública, logo esta pode ser conhecida por pessoas terceiras.
Mesmo que os sistemas assimétricos possibilitem um nível maior de
segurança, eles podem não ser adequados para criptografar grandes
volumes de dados, como ressalta \citeonline{alenezi2020}. Isso ocorre
porque a velocidade é lenta em comparação com os sistemas baseados em
chaves simétricas e eles também apresentam uma maior taxa de
utilização da \acs{CPU}. A figura abaixo ilustra o processo de criptografia
assimétrica.

\subsection{Camada 7 com confiabilidade}

Tendo em vista que, em redes, compete a camada de aplicação garantir
a segurança das informações trafegadas pelas camadas inferiores, umas
vez que esta, seja a principal no objetivo de comunicar com a outra
parte, onde é gerado todo o conteúdo que será lido pelo usuário de
destino. Por ser o ponto final de consumo e envio de dados, essa
camada se torna um alvo recorrente para ataques que buscam capturar,
manipular ou redirecionar informações sensíveis.

Em ambientes abertos e descentralizados, como é o caso de aplicações
com criptografia ponto a ponto (\acs{P2P}), os riscos são amplificados. A
segurança da aplicação não pode depender exclusivamente de mecanismos
criptográficos isolados. É fundamental que o protocolo de comunicação
utilizado seja projetado para resistir a ataques mesmo sob as piores
condições, como destaca \citeonline[p.~2, tradução
nossa]{vigano2006}, “Em redes abertas como a Internet, protocolos
devem funcionar mesmo sob suposições de pior caso, como interceptação
ou modificação de mensagens por invasores” \footnote{In open
  networks, such as the Internet, protocols should work even under
  worst-case assumptions, e.g., that messages may be eavesdropped or
tampered with by an intruder.}.

Em um cenário cada vez mais conectado, a comunicação digital segura
se tornou uma necessidade central para garantir a integridade e a
privacidade das informações trocadas entre usuários. Aplicações de
mensagens instantâneas, especialmente aquelas que operam em
estruturas descentralizadas ou ponto a ponto (\acs{P2P}), enfrentam
desafios significativos quanto à proteção de dados sensíveis.

A confidencialidade das mensagens e a identidade dos participantes
são aspectos críticos nesse contexto. A ausência de um controle
centralizador expõe os usuários a um risco maior de interceptações,
falsificações e ataques de engenharia social. Assim, a adoção de
protocolos seguros e validados é essencial para resguardar os
direitos fundamentais dos usuários, como o direito à privacidade.

Segundo \citeonline[p.~2, tradução nossa]{vigano2006}, “[...]
usuários dessas tecnologias dependem de uma infraestrutura segura
para garantir direitos e liberdades, como o direito à privacidade”
\footnote{It is also a problem for users of these technologies whose
  rights and freedoms, such as the right to privacy of personal data,
depend on a secure infrastructure.}. Essa afirmação reforça a ideia
de que a segurança da infraestrutura digital não é apenas um
requisito técnico, mas também um compromisso ético e legal, sobretudo
quando se trata da troca de mensagens pessoais ou confidenciais.

Portanto, o uso de criptografia de ponta a ponta (\textit{end-to-end
encryption}) aliada à validação formal dos protocolos utilizados se
mostra indispensável. Ferramentas como o AVISPA Tool são
relevantes neste cenário, por permitirem uma análise automatizada e
profunda de falhas em protocolos, incluindo casos em que a
criptografia permanece intacta, mas a lógica do protocolo pode ser
explorada por agentes mal-intencionados.

A proteção de dados em uma aplicação de comunicação segura, como a
proposta neste trabalho, deve ser tratada desde a concepção do
sistema, incorporando práticas como \textit{privacy by design} e
validação automatizada, de forma a garantir não apenas a
funcionalidade da comunicação, mas também a confiança de seus usuários.

\subsection{Assinatura digital}

A segurança em comunicações digitais exige que não apenas a
confidencialidade seja preservada, mas também que se garanta a
autenticidade da origem da mensagem, sua integridade, e, em
muitos casos, o não repúdio — ou seja, que o emissor não possa
negar posteriormente a autoria da mensagem enviada. Em sistemas
descentralizados ou ponto a ponto (\acs{P2P}), onde não há um servidor
central responsável por intermediar ou validar o tráfego, essa
responsabilidade recai sobre mecanismos criptográficos robustos e
cuidadosamente implementados.

A aplicação desenvolvida neste trabalho adota a arquitetura \acs{P2P} com
criptografia de ponta a ponta (\acs{E2EE}), o que garante que apenas os
dispositivos dos usuários tenham acesso ao conteúdo das mensagens.
Entretanto, a ausência de um intermediário confiável exige o uso de
assinaturas digitais como mecanismo de validação direta entre as
partes, assegurando que o conteúdo recebido não foi alterado
durante o tráfego e, mais ainda, que foi gerado pelo remetente legítimo.

Conforme apresentado por \citeonline[p.~2, tradução
nossa]{rahim2018}, “[...] as assinaturas digitais permitem ao
destinatário verificar a autenticidade e a integridade da mensagem, e
ao mesmo tempo, oferecem serviços de não repúdio”\footnote{Also,
  digital signatures also provide nonrepudiation services, which means
  protecting the sender from a claim stating that he or she has sent
information when there is no data transmission.}. Esse conceito é
essencial em ambientes onde a troca de informações sensíveis pode
acarretar consequências legais ou operacionais, e onde os servidores
não retêm cópias das mensagens para verificação posterior.

Além disso, o uso de criptografia e assinatura digital contribui para
a minimização de dados, um dos pilares da proteção de dados
moderna segundo legislações como a \acf{LGPD}. Em vez de armazenar o conteúdo das mensagens, o sistema
limita seu papel ao registro temporário da chave criptográfica
pública, garantindo que nenhuma informação sensível permaneça no
ambiente do servidor após o encerramento da sessão. Isso não só
reduz drasticamente a superfície de ataque do sistema, como
também fortalece sua conformidade com princípios legais de privacidade.

Ao incorporar práticas como assinaturas digitais baseadas em
esquemas robustos (como Ong-Schnorr-Shamir), a aplicação demonstra
que é possível manter a integridade da comunicação e responsabilizar
os envolvidos, sem comprometer a privacidade dos usuários ou criar
pontos centrais de vulnerabilidade. Trata-se de um equilíbrio
técnico e ético entre segurança, autonomia e responsabilidade, que se
alinha diretamente com as demandas atuais de sistemas distribuídos e
respeitosos à privacidade.

\subsection{\textit{Networking}}

A utilização de criptografia fim a fim (\acs{E2EE}) se destaca por sua
capacidade de eliminar o acesso do provedor ao conteúdo das
mensagens, promovendo um novo paradigma de segurança digital. Em
aplicações tradicionais, os servidores intermediários desempenham um
papel ativo no processamento das mensagens, o que os torna pontos
vulneráveis a vazamentos, acesso indevido ou mesmo requisições
judiciais de conteúdo. No entanto, a arquitetura \acs{E2EE} rompe com
essa lógica, garantindo que nem mesmo o provedor tenha acesso às
mensagens trocadas entre os usuários.

De acordo com \citeonline[p.~4, tradução nossa]{knodel2024}, “A
criptografia de ponta a ponta (\acs{E2EE}) é um padrão de segurança para
comunicação no qual apenas o remetente e o destinatário pretendido
podem ler as mensagens entre eles. Em particular, mesmo a plataforma
não consegue ler essas comunicações, apesar de ainda atuar como
intermediária”\footnote{\acfi{E2EE} is a standard
  of security for communication in which only the sender and an
  intended recipient can read communications between them. In
  particular, even the platform cannot read these communications
despite still acting as the intermediary for communication.}. Esse
modelo descentraliza o controle da informação, deslocando a
responsabilidade da proteção para as extremidades — ou seja, os
próprios dispositivos dos usuários, que são os únicos com acesso às
chaves de decriptação.

A aplicação proposta neste trabalho adota esse princípio ao extremo,
utilizando apenas as chaves criptográficas de cada usuário na
sessão, sendo excluídas ao término da comunicação e renovadas ao
início de uma nova sessão, assim como o endereço ou identificador do
usuário com o qual se deseja comunicar. O conteúdo da mensagem, por
sua vez, nunca transita em texto claro pelo servidor, nem tampouco
é armazenado. Isso garante que, mesmo em caso de invasão ou
interceptação de dados do servidor, nenhum conteúdo sensível possa
ser acessado ou reconstruído.

Essa abordagem está em conformidade com as práticas modernas de
segurança digital, que reconhecem o risco crescente do chamado
"problema do intermediário confiável". Na prática, \acs{E2EE}
transforma os servidores em meros “retransmissores cegos”,
incapazes de interpretar ou manipular o conteúdo que transportam.
Essa é uma característica fundamental para garantir não apenas a
segurança das informações, mas também a confiança do usuário no sistema.

Além disso, esse modelo fortalece o princípio do \textit{privacy by design},
ao considerar a privacidade como um requisito estrutural da
arquitetura do sistema, e não como uma camada superficial ou
opcional. O mínimo armazenamento de dados sensíveis e a exclusão
automática de informações auxiliares após a finalização da sessão de
comunicação são componentes adicionais que reforçam a
confidencialidade e reduzem a superfície de ataque do sistema.

No contexto legal, esse modelo também se alinha com os princípios
estabelecidos por regulamentações como a \acs{LGPD} (Lei nº
13.709/2018) no Brasil e o \acfi{GDPR} da União Europeia.
Ambas as legislações reforçam o conceito
de minimização de dados, exigindo que apenas os dados
estritamente necessários sejam coletados e tratados, além de garantir
transparência, finalidade específica e segurança da informação.

A criptografia \acs{E2EE}, ao garantir que os dados em trânsito sejam
inacessíveis mesmo para o controlador, mitiga riscos regulatórios e
reduz a responsabilidade legal da aplicação, principalmente no que
diz respeito à guarda e exposição de dados pessoais. Como apontado
por \citeonline[p.~4, tradução nossa]{knodel2024}, “[...]
funcionalidades de IA ou análise de
dados em sistemas \acs{E2EE} devem ser sempre opcionais e precedidas de
consentimento explícito, dada a sensibilidade e os riscos legais envolvidos”.

Nesse sentido, o sistema proposto neste trabalho também implementa
princípios de \textit{privacy by design} e \textit{privacy by default},
promovendo a privacidade como um requisito estrutural da arquitetura
e não como uma funcionalidade opcional ou complementar. Ao limitar o
armazenamento de dados sensíveis, adotar criptografia robusta, e
evitar qualquer exposição desnecessária do conteúdo ao servidor, o
sistema fortalece tanto sua segurança técnica quanto sua conformidade legal.

\subsection{Comunicação}

As comunicações digitais transformaram radicalmente a maneira como as
pessoas interagem, compartilham informações e constroem
relacionamentos. Essa transformação só foram possíveis graças ao
avanço das tecnologias que permitiram a troca de dados por meio de
dispositivos conectados. Em vez de depender exclusivamente de
encontros presenciais ou arquivos físicos, atualmente é possível
transmitir mensagens instantaneamente entre pontos distantes, de uma
forma quase imperceptível \citeonline{castells2002}.

As informações trafegam de maneira organizada, segura e eficiente,
graças a comunicação moderna. Para garantir que as mensagens cheguem
ao destino correto, os dados percorrem uma série de camadas que se
encarregam de endereçar, embalar e redirecionar os pacotes ao longo
do caminho, como afirma \citeonline{kurose2013}. Esses processos criam uma rede
no qual dispositivos, servidores e usuários trocam informações em
tempo real, de forma quase transparente ao usuário final, mas
altamente complexa em seu funcionamento interno.

Nesse contexto, \citeonline{stallings2014} aborda sobre como a preocupação com a
segurança se tornou um fator essencial por intermédio do aumento do
volume de dados e da sensibilidade das informações transmitidas,
assim como, a necessidade de garantir que apenas os envolvidos em uma
conversa tenham acesso ao conteúdo trocado. Esses métodos  de
proteção impedem o acesso indevido, mesmo em um ambiente de rede
sujeito a riscos e interferências, é possível manter a
confidencialidade e a integridade da comunicação.

Além da segurança, um aspecto igualmente essencial na comunicação
digital contemporânea é a capacidade de interatividade em tempo quase
real. A troca contínua de informações entre sistemas e usuários
permite que interações provem de maneira natural, como em conversas
instantâneas ou em ambientes colaborativos online. Para viabilizar
essa dinâmica, torna-se necessário adotar arquiteturas de comunicação
que possibilitem conexões estáveis, bidirecionais e eficientes, mesmo
diante de variações na infraestrutura da rede. Essa característica
não apenas aprimora a experiência do usuário final, como também
amplia significativamente as possibilidades de aplicação dessas
tecnologias em diversos contextos, incluindo plataformas
educacionais, serviços de atendimento ao cliente e sistemas
interativos de missão crítica \citeonline{tenenbaum2011}.

\newpage
